\documentclass[10pt,a4paper]{amsart}
\usepackage[margin=19mm]{geometry}
\usepackage{amsmath,amssymb,mathtools}
\usepackage[scr=rsfs,cal=euler]{mathalfa}
\usepackage{xfrac}
\usepackage[unicode,hidelinks,bookmarks=false]{hyperref}
\newcommand{\ie}{i.e.\ }
\newcommand{\cf}{cf.\ }
\newcommand{\resp}{respectively\ }
\newcommand{\Subsection}{\subsection}
\providecommand{\nobreakdash}{\nobreak\hbox{-}}
\setlength{\emergencystretch}{2em}
\multlinegap0pt
\usepackage{amssymb}
\newtheorem{proposition}{Proposition}
\title{Continued Fractions of Polynomial Type:\\ Theory and Encyclopedic Dictionary}
\author{Henri Cohen}
\date{Source: arXiv:2607.06581v1 (CC BY 4.0)}
\begin{document}
\maketitle

\noindent\emph{Source and licence.} Continued Fractions of Polynomial Type:\\ Theory and Encyclopedic Dictionary. Version arXiv:2607.06581v1 (CC BY 4.0), distributed by arXiv under the Creative Commons Attribution 4.0 International licence, http://creativecommons.org/licenses/by/4.0/, as recorded in the arXiv licence metadata for this exact version and shown on the abstract page https://arxiv.org/abs/2607.06581v1. Full licence notice: \texttt{licenses/arxiv-2607.06581-CC-BY-4.0.txt}.\par\smallskip

\noindent\textit{Editorial scope.} Counted unit: the article's proposition proposition (source0.tex lines 26157--26170). The article's complete original proof of it is delivered with it (source0.tex lines 26172--26231, one \texttt{proof} environment of 60 lines). No mathematics is written, repaired or re-derived: the statement and the proof are the article's own lines, in the article's own notation, and no retained line is substituted or re-flowed.  No further source-local context is needed: every cross-reference of every retained line resolves inside the two delivered spans.  Nothing was omitted from any retained span, so the package carries no omission record.  No retained line contains a citation, so the wrapper carries no bibliography and no bibliography entry.  No retained line contains a figure environment, an image-inclusion command or drawing code, so no image is delivered and the package declares no asset dependency.  Editorial formatting only, in addition: an AMS \texttt{amsart} preamble that restates the article's own declarations and macros used by the retained lines, namely the environments \texttt{proposition} on separate counters, together with the standard packages those lines need.  The retained environments print in this document as Proposition 1 in order of appearance, whereas the article prints the same items with the numbering of its own full document.  The complete native source of the article as distributed in the arXiv e-print for this exact version is bundled unchanged as \texttt{sources/204653/source0.tex}.
\begin{proposition}
  Set $(p'(n),q'(n))=(p(n+1),q(n+1))+r(n+1)(p(n),q(n))$ for
  $n\ge0$, define
  $$d(n)=r(n)(r(n+1)+a(n+1))-b(n)\;,$$
  and assume that $d(n)\ne0$ for all $n$.
  Then if we set $(p'(-1),q'(-1))=(a(0)+r(0),1)$,
  $(p'(n),q'(n))$ are the $n+1$-st convergents of the CF $(a'(n),b'(n))$
  defined by
  \begin{align*}
    a'(-1)&=a(0)+r(0)\;,\quad b'(-1)=-d(0)\;,\quad a'(0)=r(1)+a(1)\;,\\
    a'(n)&=r(n+1)+a(n+1)-r(n-1)\dfrac{d(n)}{d(n-1)}\\
    b'(n)&=b(n)\dfrac{d(n+1)}{d(n)}\;.
  \end{align*}
\end{proposition}

\begin{proof} Here and elsewhere, let $u(n)=p(n)$ or $u(n)=q(n)$.
We have
\begin{align*}
  u'(n+1)&=u(n+2)+r(n+2)u(n+1)\\
  u'(n)&=u(n+1)+r(n+1)u(n)\\
  u'(n-1)&=u(n)+r(n)u(n-1)\\
  u(n+2)&=a(n+2)u(n+1)+b(n+1)u(n)\\
  u(n+1)&=a(n+1)u(n)+b(n)u(n-1)\;,\end{align*}
so keeping $u(n)$ and $u(n-1)$ as main variables $x$ and $y$, we have
\begin{align*}
  u(n+1)&=a(n+1)x+b(n)y\\
  u(n+2)&=(a(n+2)a(n+1)+b(n+1))x+a(n+2)b(n)y
\end{align*}
Thus,
\begin{align*}
  &u'(n+1)-a'(n+1)u'(n)-b'(n)u'(n-1)\\
  &=u(n+2)+(r(n+2)-a'(n+1))u(n+1)\\
  &\phantom{=}-(a'(n+1)r(n+1)+b'(n))u(n)-b'(n)r(n)u(n-1)\\
  &=x(a(n+2)a(n+1)+b(n+1)+r(n+2)a(n+1)\\
  &\phantom{=}-a'(n+1)(a(n+1)+r(n+1))-b'(n))\\
  &\phantom{=}+y((a(n+2)+r(n+2)-a'(n+1))b(n)-b'(n)r(n))\end{align*}
giving the system
$Aa'(n+1)+Bb'(n)=U$, $Ca'(n+1)+Db'(n)=V$ with
\begin{align*}
  A&=a(n+1)+r(n+1)\;,\ B=1\;,\ C=b(n)\;,\ D=r(n)\\
  U&=(a(n+2)+r(n+2))a(n+1)+b(n+1)\\
  V&=(a(n+2)+r(n+2))b(n)\end{align*}
We will thus set
$$d(n)=AD-BC=r(n)(r(n+1)+a(n+1))-b(n)\;,$$
so in particular
$$r(n+2)+a(n+2)=(d(n+1)+b(n+1))/r(n+1)\;,$$
so
\begin{align*}
  U&=(d(n+1)+b(n+1))a(n+1)/r(n+1)+b(n+1)\\
  V&=(d(n+1)+b(n+1))b(n)/r(n+1)\;,\end{align*}
so that
\begin{align*}
  &d(n)a'(n+1)=DU-BV=r(n)U-V\\
  &=(r(n+2)+a(n+2))(r(n)a(n+1)-b(n))+r(n)b(n+2))\\
  &=(r(n+2)+a(n+2))(d(n)-r(n)r(n+1))+r(n)b(n+2)\\
  &=(r(n+2)+a(n+2))d(n)-r(n)(r(n+1)(r(n+2)+a(n+2))-b(n+2))\\
  &=(r(n+2)+a(n+2))d(n)-r(n)d(n+1)\\
  &d(n)b'(n)=AV-CU=(a(n+1)+r(n+1))V-b(n)U\\
  &=b(n)/r(n+1)((a(n+1)+r(n+1))(d(n+1)+b(n+1))\\
  &\phantom{=}-a(n+1)(d(n+1)+b(n+1))-b(n+1)r(n+1))=b(n)d(n+1)\;.\end{align*}
Thus, with $a'(n)$ and $b'(n)$ as in the proposition, for $n\ge1$ we have
$u'(n+1)=a'(n+1)u'(n)+b'(n)u'(n-1)$. Finally, we have
\begin{align*}
  (p'(-1),q'(-1))&=(a(0),1)+r(0)(1,0)=(a(0)+r(0),1)\\
  (p'(0),q'(0))&=(a(0)a(1)+b(1),a(1))+r(1)(a(0),1)\\
  &=((r(1)+a(1))a(0)+b(1),r(1)+a(1))\end{align*}
We want $p'(-1)/q'(-1)=a'(-1)$ and $q'(-1)=1$, so we must set
$p'(-1)=a'(-1)=a(0)+r(0)$. We want
$$(p'(0),q'(0))=a'(0)(p'(-1),q'(-1))+b'(-1)(1,0)=(a'(0)p'(-1)+b'(-1),a'(0)q'(-1))\;,$$
so $a'(0)=q'(0)=q(1)+r(1)q(0)=a(1)+r(1)$, hence
\begin{align*}b'(-1)&=p'(0)-a'(0)p'(-1)=p(1)+r(1)p(0)-(a(1)+r(1))(a(0)+r(0))\\
  &=a(0)a(1)+b(0)+r(1)a(0)-a(1)a(0)\\
  &\phantom{=}-r(1)a(0)-(a(1)+r(1))r(0)\\
  &=b(0)-r(0)(r(1)+a(1))=-d(0)\end{align*}
\end{proof}

\end{document}
